\section{Cuarta iteración de la Fase 3: Generador de Aleatorios}
\label{sec:GeneradorAleatorios}

Comenzamos a añadir nuevas funcionalidades al sistema del juego, y el primer
punto será añadir un generador de números aleatorios.

Los G\&W reales usaban un sistema muy sencillo. Como el contenido de las ROM
suele ser muy variado (salvo en las zonas que contienen tablas de datos), se
pueden generar números aleatorios simplemente usando un contador que apunte a
una zona de la ROM donde resida el programa, usar el valor numérico que contiene esa posición, y luego aumentar el contador. Suelen tener dos variables que marcan la zona inicial y final que se desea usar, para intentar evitar caer en la zona con tablas de
configuración, que son más homogéneas, aunque no es estrictamente necesario.

En nuestro simulador no hay ROM que leer: el cartucho es una lista de
segmentos y enlaces. El generador no puede leer memoria de programa, así
que lo sustituimos por un generador matemático que produce una secuencia
determinista a partir de una semilla: misma semilla, misma secuencia, misma
partida. Este generador alimentará cualquier punto donde se necesite generar algo que
dependa de una decisión al azar, por ejemplo las rutas de los paracaidistas, 
donde el motor pedirá \emph{dame un número de ruta} al director, y el elemento
elegido seguirá su cadena de enlaces.

\subsection{¿Dónde ubicar el generador?}

La primera decisión de esta iteración es la ubicación del generador de aleatorios. 
Hay dos candidatos evidentes: el \emph{motor de estados} o el
\emph{director del juego}.

El \textbf{motor de estados} es quien consumirá los números (elige ruta), así que
parece el punto ideal para ubicarlo, pero hay veces que el director puede cambiar la velocidad o lanzar eventos nuevos de forma aleatoria; por tanto, se ubicará en
el \textbf{director del juego}, de forma que ambos puedan usarlo si lo necesitan, siguiendo
siempre la idea de que el director es el coordinador general, y no depende del motor al que solo le lanza eventos, mientras que el motor puede consultar al director cuando lo necesite.

\subsection{El generador: LCG de 16 bits}

Se adopta un generador lineal congruencial (LCG) de 16 bits, la opción
clásica para hardware pequeño:

\begin{itemize}
	\item Estado: una sola variable de 16 bits, la \textbf{semilla}.
	
	\item Avance: $\text{estado} = (a \cdot \text{estado} + c) \pmod{2^{16}}$,
	con constantes de 16 bits (p.\,ej.,\ las clásicas $a = 25173$,
	$c = 13849$, de la familia usada por los BASIC de 8 bits). El
	desbordamiento natural del \texttt{uint16\_t} es el módulo: cero
	coste en el ESP32.
	
	\item Lectura: el número pedido entre 0 y $N-1$ se obtiene con
	$\text{estado} \bmod N$. Para los $N$ que usaremos (número de rutas,
	decenas) el sesgo del módulo es despreciable; si algún día hiciera
	falta, se anota como ampliación (mezcla de bits superiores).
\end{itemize}

\textbf{Propiedades que buscamos, y se cumplen:}

\begin{itemize}
	\item \textbf{Determinista:} misma semilla, misma secuencia, misma partida
	reproducible.
	
	\item \textbf{Sin tablas ni estado múltiple:} una variable de 16 bits y dos
	constantes; idéntico en PC y ESP32.
	
	\item \textbf{Suficientemente «aleatorio» para el juego:} no necesitamos
	calidad criptográfica, solo que las rutas no se repitan de forma
	perceptible; un LCG de 16 bits con periodo $2^{16}$ sobra para las
	decenas de tiradas de una partida.
\end{itemize}

\subsection{API y comportamiento}

En \texttt{juegoDirector} (el reloj ya reinicia la partida):

\begin{itemize}
	\item \texttt{aleatorio(max)}: avanza el LCG y devuelve un entero en el rango
	$[0, \text{max})$. Es la única puerta del generador.
	
	\item La \textbf{semilla} es una constante editable del fichero
	(fijo por código, como la \texttt{SECUENCIA} del motor), mientras no la cambiemos las partidas seguirán siempre en la misma secuencia de eventos, igual que ocurre en las máquinas reales.
	
	\item \texttt{arrancar()} reinicia el estado del LCG a la semilla, junto al
	resto del arranque (reloj, cola vacía, \texttt{iniciar()}): cada
	partida empieza con la misma secuencia.
	
	\item El diagnóstico (secuencia de números entregados) va bajo
	\texttt{GW\_DEBUG}, como en todo el proyecto.
\end{itemize}

\textbf{Cuándo se consume un número} lo indica el
ritmo de juego y se decide cuando lo requiera usar el director o el motor. Aquí solo se
construye y expone el servicio.


\subsection{Pruebas}

Para las pruebas se hace que el director inicialice el sistema, llame a una nueva función en el motor que llame 10 veces seguidas al generador de aleatorios, vuelva a inicializar y llame de nuevo a la función, verificando que ambas veces se obtienen los mismos números. Luego se elimina esa llamada en el director para poder seguir el desarrollo.


\begin{figure}[H]
	\centering
	\includegraphics[width=0.9\linewidth]{imagenes/propias/C05-I023}
	\caption[]{Pruebas del generador de números aleatorios.}
	\label{fig:c05-i023}
\end{figure}

